% Works in math mode; all special chars remain special; cheaper than \ncd.
% Will not be correct size in super and subscripts, though.
\newcommand{\ex}[1]{\mbox{\tt{#1}}}
\newcommand{\mex}[1]{\text{\ex{#1}}} % use in math mode

\newcommand{\C}{\mathbb{C}}
\newcommand{\R}{\mathbb{R}}
\newcommand{\Q}{\mathbb{Q}}
\newcommand{\Z}{\mathbb{Z}}
\newcommand{\N}{\mathbb{N}}

\newcommand{\sN}{\mathcal{N}}

\newcommand{\Pow}{\mathcal{P}}

\newcommand{\CP}{\mathbb{RP}}
\newcommand{\RP}{\mathbb{CP}}

\newcommand{\Ext}{\text{Ext}}
\newcommand{\Hom}{\text{Hom}}
\newcommand{\Tor}{\text{Tor}}

\newcommand{\compose}{\circ}

% algebra
\newcommand{\normal}{\triangleleft}
\newcommand{\isomorphic}{\cong}
\newcommand{\iso}{\cong}
\newcommand{\mono}{\hookrightarrow}
\newcommand{\epi}{\twoheadrightarrow}

% topology
\newcommand{\interior}{\text{int}}

\newcommand{\homeo}{\cong}
\newcommand{\homot}{\simeq}
\newcommand{\ra}{\rightarrow}
\newcommand{\la}{\leftarrow}

\newcommand{\intersection}{\cap}
\newcommand{\intersect}{\cap}
\newcommand{\union}{\cup}

\newcommand{\given}{\, | \, }

\newcommand{\onehalf}{\frac{1}{2}}
\newcommand{\half}{\frac{1}{2}}
\newcommand{\E}{\mathbb{E}}
\newcommand{\VAR}{\text{VAR}}
\newcommand{\COVAR}{\text{COVAR}}

\newcommand{\longline}{\text{\textemdash}}

\newcommand{\coker}{\text{coker }}
\newcommand{\im}{\text{im }}
\newcommand{\lra}[1]{\xrightarrow{#1}}
\newcommand{\lla}[1]{\xleftarrow{#1}}
\newcommand{\lda}[1]{\xdownarrow{#1}}

\newcommand{\incl}[1]{\xhookrightarrow{#1}}

\newcommand{\wedgesum}{\vee}
\newcommand{\directsum}{\oplus}
\newcommand{\tensor}{\otimes}

\newcommand{\1}{\mathbbm{1}}

\newcommand{\reduced}[1]{\widetilde{#1}}
\newcommand{\til}[1]{\widetilde{#1}}

\newcommand{\conv}{\text{conv}}
\newcommand{\cl}{\text{cl}}

\newcommand{\define}{:=}

\newcommand{\sI}{\mathcal{I}}
\newcommand{\bigunion}{\bigcup}
\newcommand{\T}[1]{\texttt{#1}}

\newcommand{\problemspace}{0.3cm}

\newcommand{\Quot}{\text{Quot}}
\newcommand{\spoly}{\text{spoly}}

\newcommand{\lcm}{\text{lcm}}
\newcommand{\LM}{\text{LM}}

\renewcommand{\phi}{\varphi}

\newcommand{\interesting}{\texttt{interesting?}}
\newcommand{\mutate}{\texttt{mutate}}

\newcommand{\env}{\Gamma}
\newcommand{\stepsTo}{\rightsquigarrow}
\newcommand{\hole}[2]{\ensuremath{\Box^{#1}_{#2}}}
\newcommand{\argshole}[1]{\ensuremath{\triangleright^{#1}}}
%{\ensuremath{\Box \hspace{-0.55em} {\hbox to 0.5em{.\hss.\hss.}}^{#1}}}
\newcommand{\targshole}[1]{\ensuremath{\triangle^{#1}}}
\newcommand{\varhole}[1]{\ensuremath{\triangle^{#1}}}
\newcommand{\generate}{\texttt{generate}}

\newcommand{\turns}{\vdash}


\newcommand{\extvar}[1]{\ensuremath{\triangleright^{#1}}}

\newcommand{\tNat}{\texttt{Nat}}
\newcommand{\tFloat}{\texttt{Float}}
\newcommand{\tList}[2]{{\texttt{List}^{#2}\;{#1}}}
\newcommand{\tProd}[3]{{\ensuremath{#1}\times^{#3}\ensuremath{#2}}}
\newcommand{\tArrow}[3]{\ensuremath{\ensuremath{#1}\ra^{#3}\ensuremath{#2}}}

%\newcommand{\ann}[1]{{\Diamond_{#1}}}
\newcommand{\ann}[1]{{#1}}


\definecolor{color-keyword}{rgb}{0.0,0.30,0.7}
\newcommand{\kw}[1]{{\text{\ttfamily\fontseries{b}\selectfont{\color{color-keyword}{#1}}}}}

\newcommand{\codearr}{
\ensuremath{
\makebox[0pt][l]{\hspace{-0.9ex}\mex{\color{color-keyword}{\kw{=}}}}
\makebox[0pt][l]{\mex{\color{color-keyword}{\kw{>}}}}
\
}
}

\newcommand{\codearri}[1]{
\ensuremath{
\makebox[0pt][l]{\hspace{-0.9ex}\mex{-}}
\makebox[0pt][l]{\mex{>}}
\
\ensuremath{{}^{#1}}
}
\!\!\!\!\!
%\code{-}
%\begin{tikzpicture}[baseline]
%    \node at (0,0) {\code{-}};
%    \node at (1.0ex,0.25ex) {\code{>}}; % Move '>' a little upwards
%\end{tikzpicture}
%\texttt{-}\raisebox{0.5ex}{\texttt{-}} \hspace{-0.15em} \raisebox{-0.5ex}{\texttt{>}}
}


\newcommand{\letexp}[3]{\mex{let\;\ensuremath{#1}\;=\;\ensuremath{#2}\;in\;\ensuremath{#3}}}
\newcommand{\lam}{\mex{$\lambda$}}
\newcommand{\lambdaexp}[2]{\mex{$\lambda$\,\ensuremath{#1}.\ensuremath{#2}}}
\newcommand{\ifexp}[3]{\mex{if\;\ensuremath{#1}\;then\;\ensuremath{#2}\;else\;\ensuremath{#3}}}
%\newcommand{\returnexp}[1]{\mex{return\;\ensuremath{#1}}}
\newcommand{\returnexp}[1]{\mex{\ensuremath{#1}}}
%\newcommand{\callexp}[2]{\mex{call\;\ensuremath{#1}\;\ensuremath{#2}}}
\newcommand{\callexp}[2]{\ensuremath{#1}\,\ensuremath{#2}}
\newcommand{\chooseexp}[2]{\callexp{\mex{choose}}{\pairexp{\ensuremath{#1}}{\ensuremath{#2}}}}
\newcommand{\chooseexpT}[3]{\callexp{\mex{choose}}{\ensuremath{#1} \; \pairexp{\ensuremath{#2}}{\ensuremath{#3}}}}

\newcommand{\withType}[2]{\ensuremath{\ensuremath{#1}\;\mex{:}\;\ensuremath{#2}}}
\newcommand{\pairexp}[2]{{\mex{(}\ensuremath{#1}\mex{,}\ensuremath{#2}\mex{)}}}

\newcommand{\acallexp}[3]{\mex{(\ensuremath{#1}\;\ensuremath{#2})${}^{\ensuremath{#3}}$}}
\newcommand{\alambdaexp}[3]{\mex{fun\;\ensuremath{#1}\;\codearri{\ensuremath{#3}}\;\ensuremath{#2}}}
\newcommand{\e}{\ensuremath{\mathit{e}}}
\newcommand{\es}{\overline{\e}}
\newcommand{\x}{\ensuremath{\mathit{x}}}
\newcommand{\xs}{\overline{\x}}
\newcommand{\type}{\ensuremath{\tau}}
\newcommand{\typep}{\ensuremath{\type'}}
\newcommand{\types}{\overline{\type}}

\newcommand{\Cxt}{\Gamma}
\newcommand{\derives}[2]{{#1\;\vdash\;#2}}

\DeclareMathOperator{\unifyOp}{unify}
\newcommand{\unify}[2]{\unifyOp(#1, #2)}
\DeclareMathOperator{\mergeOp}{merge}
\newcommand{\merge}[2]{\mergeOp(#1, #2)}
\DeclareMathOperator{\freeOp}{free}
\newcommand{\notfree}[2]{#1 \text{ not} \freeOp \text{in}\, #2}
\newcommand{\redbox}[1]{{\colorbox{red!30}{\ensuremath{#1}}}}
\newcommand{\bluebox}[1]{{\colorbox{blue!30}{\ensuremath{#1}}}}

%\newcommand{\altbox}[3]{\alt<1-#2>{\strut\ensuremath{#3}}{\colorbox{#1!30}{\strut\ensuremath{#3}}}}
\definecolor{dkred}{rgb}{0.8,0,0}
\definecolor{dkgreen}{rgb}{0,0.7,0}

\newcommand{\Checkmark}{\textcolor{dkgreen}{\checkmark}}
\newcommand{\Xmark}{\textcolor{dkred}{\usym{2613}}\,}
\newcommand{\Questionmark}{{\,\bf{?}\,}}


%\newcommand{\e}{\mex{e}}
%\newcommand{\x}{\mex{x}}
%\newcommand{\y}{\mex{y}}
%\renewcommand{\a}{\mex{a}}
%\renewcommand{\b}{\mex{b}}
%\newcommand{\z}{\mex{z}}
%\newcommand{\w}{\mex{w}}
%\newcommand{\f}{\mex{f}}
%\newcommand{\g}{\mex{g}}
%\newcommand{\h}{\mex{h}}
%\newcommand{\one}{\mex{1}}
%\newcommand{\two}{\mex{2}}

\renewcommand{\leq}{\sqsubseteq}

%\makeatletter
%\let\HL\hl
%\renewcommand\hl{%
%  \let\set@color\beamerorig@set@color
%  \let\reset@color\beamerorig@reset@color
%  \HL}
%\makeatother
%

%\newcommand\semihuge{\@setfontsize\semihuge{18.93}{23.45}}
%\newcommand\semihuge{\@setfontsize\semihuge\@xxvpt{22}}

\newcommand{\comm}[3]{\textcolor{#1}{[#2: #3]}}
\newcommand{\TODO}[1]{\comm{dkred}{TODO}{#1}}%{\colorbox{red}{TODO:#1}}

\newcommand{\rasq}{\rightsquigarrow}

\newcommand{\mybullet}{{{\color{blue}{$\blacktriangleright$}}}}

\newcommand{\unknown}{$\bullet$}
\newcommand{\lunknown}[1]{$\text{#1}_\bullet$}

\def\maxvalue(#1,#2){\ifnum #1>#2 #1\else #2\fi}
%https://tex.stackexchange.com/questions/13793/beamer-alt-command-like-visible-instead-of-like-only
\usepackage{etoolbox}
\usepackage{mathtools}

%\makeatletter
%% Detect mode. mathpalette is used to detect the used math style
%\newcommand<>\Alt[3][l]{%
%  \begingroup
%    \providetoggle{Alt@before}%
%    \alt#4{\toggletrue{Alt@before}}{\togglefalse{Alt@before}}%
%    \ifbool{mmode}{%
%      \expandafter\mathpalette
%      \expandafter\math@Alt
%    }{%
%      \expandafter\make@Alt
%    }%
%    {{#1}{#2}{#3}}%
%  \endgroup
%}
%
%% Un-brace the second argument (required because \mathpalette reads the three arguments as one
%\newcommand\math@Alt[2]{\math@@Alt{#1}#2}
%
%% Set the two arguments in boxes. The math style is given by #1. \m@th sets \mathsurround to 0.
%\newcommand\math@@Alt[4]{%
%  \setbox\z@ \hbox{$\m@th #1{#3}$}%
%  \setbox\@ne\hbox{$\m@th #1{#4}$}%
%  \@Alt{#2}%
%}
%
%% Un-brace the argument
%\newcommand\make@Alt[1]{\make@@Alt#1}
%
%% Set the two arguments into normal boxes
%\newcommand\make@@Alt[3]{%
%  \sbox\z@ {#2}%
%  \sbox\@ne{#3}%
%  \@Alt{#1}%
%}
%
%% Place one of the two boxes using \rlap and place a \phantom box with the maximum of the two boxes
%\newcommand\@Alt[1]{%
%  \setbox\tw@\null
%  \ht\tw@\ifnum\ht\z@>\ht\@ne\ht\z@\else\ht\@ne\fi
%  \dp\tw@\ifnum\dp\z@>\dp\@ne\dp\z@\else\dp\@ne\fi
%  \wd\tw@\ifnum\wd\z@>\wd\@ne\dimexpr\wd\z@/2\relax\else\dimexpr\wd\@ne/2\relax\fi
%  %
%  \ifstrequal{#1}{l}{%
%    \rlap{\iftoggle{Alt@before}{\usebox\z@}{\usebox\@ne}}%
%    \copy\tw@
%    \box\tw@
%  }{%
%    \ifstrequal{#1}{c}{%
%      \copy\tw@
%      \clap{\iftoggle{Alt@before}{\usebox\z@}{\usebox\@ne}}%
%      \box\tw@
%    }{%
%      \ifstrequal{#1}{r}{%
%        \copy\tw@
%        \box\tw@
%        \llap{\iftoggle{Alt@before}{\usebox\z@}{\usebox\@ne}}%
%      }{%
%      }%
%    }%
%  }%
%}
%%\makeatother
%\newcommand<>\Alt[2]{{%
%    \sbox0{$\displaystyle #1$}%
%    \sbox1{$\displaystyle #2$}%
%    \alt#3%
%        {\parbox[c][\maxvalue(\ht0,\ht1)][c]{\maxvalue(\wd0,\wd1)}
%          {\centering\usebox0}}
%        {\parbox[c][\maxvalue(\ht0,\ht1)][c]{\maxvalue(\wd0,\wd1)}
%          {\centering\usebox1}}
%        %{\rlap{\usebox0}\vphantom{\usebox1}\hphantom{\ifnum\wd0>\wd1 \usebox0\else\usebox1\fi}}%
%        %{\rlap{\usebox1}\vphantom{\usebox0}\hphantom{\ifnum\wd0>\wd1 \usebox0\else\usebox1\fi}}%
%}}
%
%\newcommand{\highlight}[2]{\colorbox{#1!30}{\strut\ensuremath{#2}}}
%
%\newcommand<>\highlightalt[2]{\alt#3{\colorbox{#1!30}{\strut\ensuremath{#2}}}
%                                    {\vphantom{\colorbox{#1!30}{\strut\ensuremath{#2}}}
%                                     \makebox[\widthof{\colorbox{#1!30}{\strut\ensuremath{#2}}}]{
%                                     \strut\ensuremath{#2}
%                                     }}
%                                    }
%\newcommand<>\underlinealt[1]{\Alt#2{\fbox{\ensuremath{#1}}}{#1}}
\definecolor{dkred}{rgb}{0.8,0,0}
\definecolor{dkgreen}{rgb}{0,0.7,0}
\definecolor{dkblue}{rgb}{0,0,0.7}
\definecolor{dkyellow}{rgb}{0.8,0.8,0}
\definecolor{dkpurple}{rgb}{0.7,0,0.7}
\definecolor{dkgray}{rgb}{0.8,0.8,0.8}
\definecolor{dkorange}{rgb}{0.85,0.33,0.0}

\renewcommand{\leq}{\sqsubseteq}
\newcommand{\join}{\sqcup}

\newcommand{\wrap}{\mathit{wrap}}
\newcommand{\unwrap}{\mathit{unwrap}}

%\newcommand{\altbox}[3]{\alt<1-#2>{\vphantom{\colorbox{#1!30}{\strut\ensuremath{#3}}}
%                                   \makebox[\widthof{\colorbox{#1!30}{\strut\ensuremath{#3}}}]{
%                                   \strut\ensuremath{#3}
%                                   }}
%                                  {\colorbox{#1!30}{\strut\ensuremath{#3}}}}
%\newcommand{\altboxm}[3]{\alt<1-#2>{\makebox[0pt][l]{\strut #3}\phantom{\colorbox{#1!30}{\strut #3}}}
%                                  {\colorbox{#1!30}{\strut #3}}}
%
%\newcommand{\altboxa}[3]{
%  \alt<#2>
%       {\colorbox{#1!30}{\strut #3}}
%       {\makebox[0pt][l]{\strut #3}\phantom{\colorbox{#1!30}{\strut #3}}}
%}

\newcommand{\altcolor}[3]{%
  \alt<#1>{%
    \tikz[baseline]{%
      \node[
        fill={#2},
        fill opacity=0.3,
        text opacity=1,
        anchor=base,
        rounded corners=1pt,
        inner xsep=1pt,
        inner ysep=1pt
      ] {\strut #3};%
    }%
  }{%
    \tikz[baseline]{%
      \node[
        fill=none,
        text opacity=1,
        anchor=base,
        inner xsep=1pt,
        inner ysep=1pt
      ] {\strut #3};%
    }%
  }%
}

\newcommand{\altborder}[3]{%
  \alt<#1>{%
    \tikz[baseline]{%
      \node[
        draw={#2},
        draw opacity=0.5,
        very thick,
        fill=none,
        text opacity=1,
        anchor=base,
        rounded corners=1pt,
        inner xsep=1pt,
        inner ysep=1pt
      ] {\strut #3};%
    }%
  }{%
    \tikz[baseline]{%
      \node[
        draw=none,
        fill=none,
        text opacity=1,
        anchor=base,
        inner xsep=1pt,
        inner ysep=1pt
      ] {\strut #3};%
    }%
  }%
}

\newcommand{\altunderline}[2]{%
  \alt<#1>{%
    \tikz[baseline]{%
      \node[
        draw=none,
        very thick,
        fill=none,
        text opacity=1,
        anchor=base,
        rounded corners=1pt,
        inner xsep=1pt,
        inner ysep=1pt,
        path picture={
          \draw[black, very thick, opacity=0.5] (current bounding box.south west)
                         -- (current bounding box.south east);  % Bottom border only
        }
      ] {\strut {#2}};%
    }%
  }{%
    \tikz[baseline]{%
      \node[
        draw=none,
        fill=none,
        text opacity=1,
        anchor=base,
        inner xsep=1pt,
        inner ysep=1pt
      ] {\strut #2};%
    }%
  }%
}

\definecolor{dkyelloworange}{rgb}{0.9,0.7,0.2}
\definecolor{dklightblue}{rgb}{0.6,0.6,1.0}
\definecolor{dkgreen}{rgb}{0.4,0.8,0.4}
\definecolor{dkdarkgray}{rgb}{0.1,0.1,0.1}
\definecolor{dkdarkgreen}{rgb}{0.1,0.5,0.1}
\newcommand{\webid}{
  \tikz[baseline]{%
      \node[
        fill={dkyelloworange},
        fill opacity=0.3,
        text opacity=1,
        anchor=base,
        rounded corners=1pt,
        inner xsep=1pt,
        inner ysep=3pt
      ] {\ensuremath{u}};%
  }
}
%\colorbox{pink!70}{\ensuremath{u}}}
\newcommand{\webidp}{
  \tikz[baseline]{%
      \node[
        fill={dkdarkgreen},
        fill opacity=0.3,
        text opacity=1,
        anchor=base,
        rounded corners=1pt,
        inner xsep=1pt,
        inner ysep=3pt
      ] {\ensuremath{v}};%
  }
}
%\colorbox{blue!30}{\ensuremath{v}}}

\newcommand{\gray}[1]{
  {\color{dkgray}#1}
}

\newcommand{\tightldots}{.\hspace{-0.25em}.\hspace{-0.25em}.\hspace{-0.15em}}

\newcommand{\codequote}{{\hspace{-0.1em}\textnormal{'}}}

\newcommand<>\Alt[2]{{%
    \sbox0{$\displaystyle #1$}%
    \sbox1{$\displaystyle #2$}%
    \alt#3%
        {\parbox[c][\maxvalue(\ht0,\ht1)][c]{\maxvalue(\wd0,\wd1)}
          {\centering\usebox0}}
        {\parbox[c][\maxvalue(\ht0,\ht1)][c]{\maxvalue(\wd0,\wd1)}
          {\centering\usebox1}}
        %{\rlap{\usebox0}\vphantom{\usebox1}\hphantom{\ifnum\wd0>\wd1 \usebox0\else\usebox1\fi}}%
        %{\rlap{\usebox1}\vphantom{\usebox0}\hphantom{\ifnum\wd0>\wd1 \usebox0\else\usebox1\fi}}%
}}
\newcommand<>\boxalt[1]{\Alt#2{\fbox{\ensuremath{#1}}}{#1}}


% https://tex.stackexchange.com/questions/85269/easy-absolute-positioning-in-beamer
\usepackage{etoolbox}

\newcount\mycount

\def\mytempcoord{}
\def\mytempopa{}
%
\newcommand\doatpos[3]{%
    \bgroup
    \renewcommand*{\do}[1]{\ifnum\mycount<2
                        \edef\mytempcoord{\mytempcoord ##1%
                            \ifnum\mycount<1
                            ,
                            \fi
                        }
                        \else
                        \edef\mytempopa{##1}
                        \fi
                        \advance\mycount by 1
                        }
    \docsvlist{#2}\global\advance\mycount by -3
    \begin{tikzpicture}[remember picture, overlay]
    \node[anchor=#1, opacity=\mytempopa] at (\mytempcoord) {#3};
    \end{tikzpicture}%
        %
    %
    \egroup
%
}

\definecolor{dklabel}{rgb}{0.0,0.0,0.0}
\newcommand{\lambdalabel}[1]{\ensuremath{{\color{dklabel}{\ell_{#1}}}}}
\newcommand{\calllabel}[1]{\ensuremath{{\color{dklabel}{c_{#1}}}}}
\newcommand{\unknownlambdalabel}{\ensuremath{{\color{dklabel}{\ell_{\bullet}}}}}
\newcommand{\unknowncalllabel}{\ensuremath{{\color{dklabel}{c_{\bullet}}}}}

\definecolor{dkcomment}{rgb}{0.85,0.33,0.0}

\newcommand{\LargeHuge}{\fontsize{18}{21}\selectfont}
